\documentclass[12pt,oneside]{book}

% ============================================================
% METADATOS
% ============================================================
\newcommand{\universidad}{Universidad de Buenos Aires}
\newcommand{\facultad}{Facultad de Derecho}
\newcommand{\materia}{Derechos Reales}
\newcommand{\titulo}{Derecho de Dominio y sus L\'imites}
\newcommand{\autor}{Isaac Edgar Camacho Ocampo}
\newcommand{\anio}{2025}

% ============================================================
% PAQUETES
% ============================================================
\usepackage[utf8]{inputenc}
\usepackage[T1]{fontenc}
\usepackage[spanish,es-nodecimaldot]{babel}

\usepackage[
    top=2.5cm,
    bottom=2.5cm,
    left=3cm,
    right=2.5cm,
    headheight=15pt
]{geometry}

\usepackage{amsmath}
\usepackage{amssymb}
\usepackage{mathtools}

\usepackage{graphicx}
\usepackage{float}
\usepackage{caption}
\usepackage{subcaption}

\usepackage{booktabs}
\usepackage{array}
\usepackage{longtable}
\usepackage{tabularx}

\usepackage{enumitem}

\usepackage{xcolor}
\usepackage[most]{tcolorbox}

\usepackage{fancyhdr}
\usepackage{ragged2e}

\usepackage[
    colorlinks=true,
    linkcolor=blue!50!black,
    citecolor=green!40!black,
    urlcolor=blue!60!black,
    pdftitle={\titulo},
    pdfauthor={\autor},
    pdfsubject={Apuntes universitarios},
    bookmarks=true
]{hyperref}

% ============================================================
% CONFIGURACIÓN
% ============================================================
\graphicspath{
    {./img/}
    {./graficos/}
    {./figuras/}
}

\setcounter{tocdepth}{2}

\setlength{\parindent}{0pt}
\setlength{\parskip}{0.5em}

\setlist[itemize]{
    topsep=0.3em,
    itemsep=0.2em
}

\setlist[enumerate]{
    topsep=0.3em,
    itemsep=0.2em
}

\clubpenalty=10000
\widowpenalty=10000

% ============================================================
% COMANDOS
% ============================================================
\newcommand{\abs}[1]{\left|#1\right|}
\newcommand{\norm}[1]{\left\lVert#1\right\rVert}
\newcommand{\vect}[1]{\mathbf{#1}}
\newcommand{\deriv}[2]{\frac{d#1}{d#2}}
\newcommand{\pderiv}[2]{\frac{\partial#1}{\partial#2}}

% ============================================================
% ENTORNOS / CAJAS
% ============================================================
\newtcolorbox{definition}[1]{
    enhanced,
    breakable,
    title={Definición: #1},
    fonttitle=\bfseries,
    colback=blue!3,
    colframe=blue!50!black,
    boxrule=0.8pt,
    arc=2mm
}

\newtcolorbox{important}{
    enhanced,
    breakable,
    title={Importante},
    fonttitle=\bfseries,
    colback=yellow!8,
    colframe=orange!70!black,
    boxrule=0.8pt,
    arc=2mm
}

\newtcolorbox{example}[1]{
    enhanced,
    breakable,
    title={Ejemplo: #1},
    fonttitle=\bfseries,
    colback=green!4,
    colframe=green!40!black,
    boxrule=0.8pt,
    arc=2mm
}

\newtcolorbox{formula}[1]{
    enhanced,
    breakable,
    title={Fórmula / Regla: #1},
    fonttitle=\bfseries,
    colback=purple!4,
    colframe=purple!50!black,
    boxrule=0.8pt,
    arc=2mm
}

\newtcolorbox{exercise}[1]{
    enhanced,
    breakable,
    title={Ejercicio: #1},
    fonttitle=\bfseries,
    colback=gray!5,
    colframe=gray!60!black,
    boxrule=0.8pt,
    arc=2mm
}

\newtcolorbox{note}{
    enhanced,
    breakable,
    title={Nota},
    fonttitle=\bfseries,
    colback=white,
    colframe=gray!50,
    boxrule=0.6pt,
    arc=2mm
}

% ============================================================
% CONFIGURACIÓN FINAL
% ============================================================
\pagestyle{fancy}
\fancyhf{}

\fancyhead[L]{\nouppercase{\leftmark}}
\fancyhead[R]{\materia}

\fancyfoot[C]{\thepage}

\renewcommand{\headrulewidth}{0.4pt}
\renewcommand{\footrulewidth}{0pt}

\fancypagestyle{plain}{
    \fancyhf{}
    \fancyfoot[C]{\thepage}
    \renewcommand{\headrulewidth}{0pt}
}

\renewcommand{\arraystretch}{1.25}

% ============================================================
% DOCUMENTO
% ============================================================
\begin{document}

% ------------------------------------------------------------
% PORTADA
% ------------------------------------------------------------
\begin{titlepage}

\thispagestyle{empty}

\begin{center}

\vspace*{1cm}

{\large \textsc{\universidad}}

\vspace{0.2cm}

{\large \textsc{\facultad}}

\vspace{3cm}

{\Huge \textbf{\materia}}

\vspace{1cm}

{\LARGE \titulo}

\vspace{0.8cm}

\textit{Comprender los conceptos es el primer paso para convertir el estudio en conocimiento.}

\vspace{1.2cm}

\rule{0.70\textwidth}{0.5pt}

\vspace{0.5cm}

{\large Apuntes de estudio}

\vspace{0.5cm}

\rule{0.70\textwidth}{0.5pt}

\vfill

{\large \autor}

\vspace{0.3cm}

{\large \anio}

\end{center}

\vfill

\begin{flushleft}
\textit{Este es un modesto aporte a los estudiantes de ``\materia'' no pretende sustituir el material oficial de la materia.}
\end{flushleft}

\end{titlepage}

% ------------------------------------------------------------
% ÍNDICE
% ------------------------------------------------------------
\tableofcontents

% ------------------------------------------------------------
% CONTENIDO
% ------------------------------------------------------------

% ============================================================
\chapter[El derecho de dominio]{El derecho de dominio}
% ============================================================

\section{Panorama de los temas}

Quien es dueño de una casa puede, como propietario, hacer casi todo con ella: venderla, alquilarla, reformarla. Esa propiedad puede ser plena o estar afectada por una hipoteca o un usufructo. Cabe preguntarse, además, cómo se llegó a ser propietario: por compra con escritura o por herencia. Y, aun siendo dueño, no es posible hacer ruido hasta las 3 de la mañana, tapar la ventana del vecino ni bloquear el río. Ese es el recorrido de estos apuntes: concepto y caracteres del dominio, su clasificación, su adquisición y, finalmente, los límites a los que está sometido.

\section{Las disposiciones comunes aplicadas al dominio}

Cada vez que se aborda un derecho real en particular, resulta metodológicamente fundamental \textbf{exponer en acto} todo lo aprendido sobre las disposiciones comunes a los derechos reales. Este recurso es sumamente importante para el estudio, para el examen parcial y para el ejercicio profesional, ya que implica aplicar la nueva metodología del Código Civil y Comercial (CCyC) en la actividad cotidiana.

El CCyC, además de los principios rectores aplicables a todo el ordenamiento (buena fe, previsión del abuso del derecho, de la posición dominante y preeminencia del orden público), utiliza partes generales y disposiciones comunes. Estas disposiciones comunes no se repiten en la regulación de cada derecho real en particular.

Aplicado al dominio, el ejercicio consiste en afirmar de él lo siguiente:

\begin{itemize}
    \item Es un \textbf{derecho real} y, por tanto, un poder jurídico.
    \item Es de \textbf{estructura legal}, porque todos los derechos reales lo son.
    \item Se ejerce directamente sobre un objeto y, muy excepcionalmente, sobre un derecho cuando la ley específicamente lo indique.
    \item Sus normas son de \textbf{orden público}.
    \item Es un derecho real sobre \textbf{cosa propia}, que se ejerce por la \textbf{posesión}.
    \item Se adquiere con \textbf{título suficiente y modo}.
    \item Puede adquirirse en forma \textbf{originaria o derivada}.
    \item Las normas que lo regulan no pueden ser modificadas por la voluntad de los particulares, bajo pena de nulidad.
\end{itemize}

\section{Dominio y propiedad}

Para determinar si «dominio» y «propiedad» son sinónimos, y qué protege la Constitución Nacional (CN) cuando garantiza el derecho de propiedad, es necesario distinguir dos acepciones del término «propiedad».

\begin{table}[H]
\centering
\caption{Propiedad en sentido amplio y en sentido restringido}
\small
\begin{tabularx}{\textwidth}{
    >{\raggedright\arraybackslash}p{0.22\textwidth}
    >{\raggedright\arraybackslash}X
    >{\raggedright\arraybackslash}p{0.17\textwidth}
}
\toprule
\textbf{Concepto} & \textbf{Alcance} & \textbf{Referencia normativa} \\
\midrule
Propiedad en sentido amplio &
Abarca tanto derechos personales como derechos reales (bienes materiales e inmateriales que integran el patrimonio de una persona). &
CN, arts.~14 y 17 \\
Propiedad en sentido restringido (técnico) &
Es el derecho real de dominio: el más extenso y completo que puede recaer sobre una cosa determinada. &
CCyC, art.~1941 \\
\bottomrule
\end{tabularx}
\end{table}

\begin{itemize}
    \item \textbf{CN, art.~14:} garantiza a todos los habitantes de la Nación usar y disponer de la propiedad.
    \item \textbf{CN, art.~17:} la propiedad es inviolable.
\end{itemize}

La Constitución Nacional se refiere a la propiedad en sentido amplio: protege tanto derechos personales como derechos reales, dado que en su conjunto integran el patrimonio de una persona.

\section{Regulación normativa del dominio}

El derecho de dominio está regulado en tres niveles normativos:

\begin{itemize}
    \item Título Tercero del Libro Cuarto del CCyC, arts.~1941 a 1982 (regulación específica).
    \item Disposiciones generales en materia de derechos reales, arts.~1882 a 1907.
    \item Principios generales del ordenamiento civil y comercial: arts.~10, 11, 12 y 14.
\end{itemize}

\section{Concepto de dominio}

\begin{definition}{Dominio perfecto (art.~1941 CCyC)}
El dominio es el derecho real en virtud del cual una cosa se encuentra sometida a la voluntad y a la acción de una persona. El dominio pleno o perfecto es el que reúne las facultades de usar, gozar y disponer material y jurídicamente de una cosa, dentro de los límites previstos por la ley. Se presume perfecto hasta que se pruebe lo contrario.
\end{definition}

El dominio es el derecho más extenso que puede tener una persona sobre una cosa: confiere la facultad de \textbf{usar, gozar y disponer} material y jurídicamente sobre ella. Los demás derechos reales no son sino fragmentos del dominio: son \textbf{desmembraciones} que el propio titular, que reúne en sí mismo todas esas facultades, va constituyendo al separar alguna de ellas y conferirla a otro, ya sea mediante derechos de disfrute o mediante derechos de garantía.

\begin{important}
Cuando se habla de «dominio» sin más, se hace referencia al dominio perfecto. Esto es así porque el art.~1941 establece que el dominio se presume perfecto hasta que se pruebe lo contrario.
\end{important}

\section{Caracteres del dominio perfecto}

Los tres caracteres del dominio perfecto son: \textbf{absoluto}, \textbf{exclusivo} y \textbf{perpetuo}.

\subsection{Carácter absoluto}

El dominio es \textbf{absoluto} porque otorga todas las facultades ---usar, gozar y disponer, material y jurídicamente--- reunidas en una misma persona. No existe ningún otro titular de derecho real que reúna sobre sí semejante cantidad de facultades.

Que el dominio sea absoluto no significa que el titular pueda hacer «lo que quiere» en sentido irrestricto. El ejercicio del dominio no debe ser abusivo y se encuentra limitado por las normas de orden público, ya sea en interés de la comunidad o de la vecindad (véase el capítulo sobre los límites al dominio).

\subsection{Carácter exclusivo}

\textbf{Art.~1943 CCyC:} el dominio es exclusivo y no puede tener más de un titular.

La exclusividad es un carácter esencial: sin exclusividad no hay dominio. Si dos personas tuvieran sobre la misma cosa un derecho real, ya no sería dominio sino \textbf{condominio}, que es un derecho rotundamente diferente, pues ha aparecido otro condueño con características absolutamente distintas.

\textbf{Art.~1944 CCyC (facultad excluyente):} el dominio es excluyente; el propietario puede excluir a cualquiera del uso y goce de la cosa, remover los objetos puestos en ella y encerrar los inmuebles con muros, cercos o fosos.

\subsection{Carácter perpetuo}

El dominio no tiene límites en el tiempo y subsiste con independencia de que el titular lo ejerza o no. Se perpetúa incluso más allá de la vida humana, razón por la cual existe el sistema de sucesiones.

Todos los demás deben respetar el vínculo del titular de dominio con la cosa. Quien, por su parte, poseyera la cosa por el tiempo y con las calidades requeridas por la ley, adquirirá el dominio por \textbf{prescripción adquisitiva}.
% TODO: contenido incompleto o ambiguo en las notas originales (la formulación original de este párrafo es confusa).

\section{Facultades del titular de dominio}

Las facultades del titular de dominio son sumamente amplias:

\begin{itemize}
    \item \textbf{Disponer jurídicamente:} incluye \textbf{enajenar} (transmitir) a título oneroso (venta, permuta) o a título gratuito (donación, dación en pago). El término «enajenar» no es sinónimo de vender, sino de transmitir.
    \item \textbf{Constituir derechos reales sobre la cosa:} por ejemplo, usufructo, hipoteca, servidumbre.
    \item \textbf{Administrar:} por ejemplo, dar la cosa en comodato (gratuito u oneroso).
    \item \textbf{Abandonar:} cuando el abandono recae sobre inmuebles, debe hacerse mediante escritura pública.
    \item \textbf{Disponer materialmente:} construir, destruir (con los límites legales) y transformar la cosa.
    \item \textbf{Usar y gozar:} obtener los frutos (civiles: locaciones; naturales: cosechas, etc.).
\end{itemize}

\begin{example}{El usufructo}
El titular de dominio tiene todas las facultades. Cuando constituye un usufructo, «saca» de lo que tiene ---el uso y el goce--- y los pone en el usufructuario. Este será un titular de derecho real que no es propietario de la cosa y que, sin embargo, tiene la fortaleza de los derechos reales: orden público, publicidad, contenido patrimonial y protección legal. El usufructuario tiene todo lo que tiene el titular de dominio, menos la propiedad.
\end{example}

\section{Extensión del dominio}

\textbf{Art.~1945 CCyC (extensión):} el dominio de una cosa comprende los objetos que forman un todo con ella o son sus accesorios. El dominio de un inmueble se extiende al subsuelo y al espacio aéreo. Se presume que pertenecen al dueño del inmueble todas las construcciones, plantaciones y forestaciones existentes.

Como decían los romanos, el dueño de un inmueble lo es «desde el cielo hasta el infierno». La extensión del dominio abarca:

\begin{itemize}
    \item Todos los accesorios de la cosa (por ejemplo, las cañerías de agua y gas dentro de las paredes de un inmueble; las ruedas de un vehículo).
    \item El subsuelo del inmueble.
    \item El espacio aéreo sobre el inmueble.
    \item Todas las construcciones, forestaciones y plantaciones (salvo propiedad horizontal o derecho de superficie).
\end{itemize}

\begin{example}{Accesorios y construcciones}
Quien compra un inmueble adquiere junto con él las cañerías de agua y gas incrustadas en sus paredes. Quien compra un vehículo adquiere también las ruedas, aunque individualmente sean cosas muebles cuando están separadas. En la práctica profesional debe tenerse presente que todo lo que está construido en el terreno pertenece al dueño del terreno.
\end{example}

\section{Clasificación del dominio}

El art.~1946 establece que los tipos de dominio imperfecto son únicamente tres y que son de orden público.

\begin{table}[H]
\centering
\caption{Tipos de dominio (art.~1946 CCyC)}
\small
\begin{tabularx}{\textwidth}{
    >{\raggedright\arraybackslash}p{0.17\textwidth}
    >{\raggedright\arraybackslash}X
    >{\raggedright\arraybackslash}p{0.14\textwidth}
    >{\raggedright\arraybackslash}p{0.17\textwidth}
}
\toprule
\textbf{Tipo} & \textbf{Característica principal} & \textbf{Afecta al carácter} & \textbf{Normativa} \\
\midrule
Dominio perfecto & Sin gravamen, sin plazo, sin condición & Ninguno & Art.~1941 \\
Dominio revocable & Sometido a plazo o condición resolutoria (máximo 10 años) & Perpetuo & Arts.~1946 y 1965 \\
Dominio fiduciario & Ejercido durante un plazo para luego transmitir & Perpetuo & Arts.~1701 y ss. \\
Dominio desmembrado & Gravado con derecho real a favor de tercero & Absoluto & Normas propias de cada derecho \\
\bottomrule
\end{tabularx}
\end{table}

\subsection{Dominio revocable}

\textbf{Art.~1965 CCyC:} el dominio revocable es aquel sometido a condición o plazo resolutorios. El plazo máximo de la condición resolutoria no puede superar los 10 años. Cumplida la condición o vencido el plazo, el dueño debe transmitir la cosa a quien se la transmitió originariamente.

Si al cabo de diez años no se revoca el dominio, este queda convertido en dominio perfecto.

\subsection{Dominio fiduciario}

\textbf{Arts.~1701 y ss. CCyC (fideicomiso):} el dominio fiduciario está regulado en el Libro Tercero, conjuntamente con el contrato de fideicomiso, y no en el Libro de Derechos Reales. Es ejercido por el adquirente durante un determinado tiempo, al cabo del cual deberá necesariamente transmitir la propiedad.

\begin{important}
Debe distinguirse: el \textbf{fideicomiso} es el contrato; el \textbf{dominio fiduciario} es el derecho real.
\end{important}

\subsection{Dominio desmembrado}

El dominio desmembrado es el que se encuentra gravado con un derecho real a favor de un tercero (por ejemplo, usufructo o hipoteca). Se rige por las normas propias de cada derecho real constituido y afecta el carácter absoluto del dominio.

\section{Cuadro Cornell: concepto y caracteres del dominio}

\begin{longtable}{
    >{\raggedright\arraybackslash}p{0.26\textwidth}
    >{\raggedright\arraybackslash}p{0.66\textwidth}
}
\toprule
\textbf{Preguntas / palabras clave} & \textbf{Notas / respuestas} \\
\midrule
\endfirsthead
\toprule
\textbf{Preguntas / palabras clave} & \textbf{Notas / respuestas} \\
\midrule
\endhead
\midrule
\multicolumn{2}{r}{\small\itshape continúa en la página siguiente} \\
\endfoot
\bottomrule
\endlastfoot

\textbf{¿Son sinónimos «dominio» y «propiedad»?} &
No. «Propiedad» en sentido amplio (CN, arts.~14 y 17) abarca derechos personales y reales. «Dominio» es la propiedad en sentido técnico: el derecho real más extenso y completo sobre una cosa determinada (art.~1941 CCyC). \\

\textbf{¿Qué es el dominio según el art.~1941?} &
El derecho real en virtud del cual una cosa se encuentra sometida a la voluntad y a la acción de una persona. El dominio perfecto reúne las facultades de usar, gozar y disponer material y jurídicamente de la cosa, dentro de los límites legales, y se presume perfecto hasta prueba en contrario. \\

\textbf{¿Cuáles son los caracteres del dominio perfecto?} &
Absoluto (todas las facultades reunidas en un titular), exclusivo (un solo titular, art.~1943) y perpetuo (sin límite temporal, independiente de su ejercicio). Si alguno falta o se afecta, ya no hay dominio perfecto. \\

\textbf{¿Qué significa que el dominio es absoluto?} &
Que reúne las máximas facultades: usar, gozar y disponer material y jurídicamente. No significa que sea ilimitado: el ejercicio no puede ser abusivo y está acotado por las normas de orden público. \\

\textbf{¿Qué ocurre si dos personas tienen un derecho real sobre la misma cosa?} &
Ya no hay dominio sino condominio, porque la exclusividad es un carácter esencial del dominio. \\

\textbf{¿Qué abarca la extensión del dominio (art.~1945)?} &
Los objetos que forman un todo con la cosa o son sus accesorios, el subsuelo, el espacio aéreo y todas las construcciones, plantaciones y forestaciones existentes (salvo propiedad horizontal o derecho de superficie). \\

\textbf{¿Cuáles son los tipos de dominio imperfecto?} &
El art.~1946 CCyC establece tres, solo tres, de orden público: (1) revocable, sujeto a plazo o condición resolutoria, máximo 10 años; (2) fiduciario, en el marco del fideicomiso (arts.~1701 y ss.); (3) desmembrado, gravado con un derecho real a favor de tercero (usufructo, hipoteca, etc.). \\

\midrule
\multicolumn{2}{p{0.92\textwidth}}{
\textbf{Resumen:}
El dominio es el más extenso de los derechos reales: permite usar, gozar y disponer material y jurídicamente de la cosa, y se presume perfecto (absoluto, exclusivo y perpetuo) hasta que se demuestre lo contrario. Los demás derechos reales son desmembraciones de sus facultades. El CCyC regula tres tipos de dominio imperfecto ---revocable, fiduciario y desmembrado---, todos de orden público.
} \\
\end{longtable}

% ============================================================
\chapter[Adquisición del dominio]{Adquisición del dominio}
% ============================================================

\section{Modos de adquisición}

Los modos de adquirir el dominio son aquellos hechos y actos jurídicos que la ley prevé y admite para adjudicar en propiedad una cosa mueble o un inmueble a una persona. Los hechos ilícitos (hurto, robo) no son modos reconocidos.

\textbf{Arts.~1947 a 1963 CCyC:} regulan los distintos modos de adquisición del dominio, incluyendo una casuística específica: animales de caza, enjambres, pesca, acrecimiento de inmuebles por hechos de la naturaleza (aluvión, avulsión), plantaciones, etc.

\begin{table}[H]
\centering
\caption{Clasificación de los modos de adquisición del dominio}
\small
\begin{tabularx}{\textwidth}{
    >{\raggedright\arraybackslash}p{0.15\textwidth}
    >{\raggedright\arraybackslash}p{0.16\textwidth}
    >{\raggedright\arraybackslash}X
    >{\raggedright\arraybackslash}p{0.22\textwidth}
}
\toprule
\textbf{Clasificación} & \textbf{Subtipo} & \textbf{Descripción} & \textbf{Ejemplo} \\
\midrule
Por el acto & Entre vivos & Transmisión en vida de ambas partes & Compraventa, donación \\
Por el acto & Mortis causa & Transmisión por fallecimiento & Herencia \\
Por el objeto & A título universal & Se adquieren todos los derechos del causante & Herencia universal \\
Por el objeto & A título singular & Se adquiere un bien determinado & Legado de cosa cierta \\
Por la causa & Originario & La ley permite tomar la cosa por sí mismo & Apropiación, prescripción \\
Por la causa & Derivado & Se transmite de un titular anterior & Tradición (título + modo) \\
\bottomrule
\end{tabularx}
\end{table}

\subsection{Apropiación (modo originario)}

\textbf{Art.~1947 CCyC:} son susceptibles de apropiación las cosas muebles sin dueño o abandonadas por su dueño. Los animales domésticos y los animales domesticados no son apropiables.

\subsection{Modos derivados: título y modo}

Los modos derivados de adquisición remiten a la teoría del \textbf{título y el modo}. El modo derivado por excelencia es la \textbf{tradición}. Para adquirir el dominio por esta vía se necesitan dos elementos: \textbf{título suficiente} (acto jurídico válido que justifica la adquisición) y \textbf{modo} (tradición: entrega efectiva de la cosa). Uno sin el otro no es suficiente.

\begin{important}
Si un poseedor (no titular) transmite la posesión, solo transmite el hecho de la posesión. Pero si el propietario (legitimado, con capacidad y título suficiente) transmite, lo que transmite es el derecho de poseer.
\end{important}

\section{Protección del sub adquirente de buena fe a título oneroso}

Este tema refleja la fortaleza del derecho real: ante diversas personas con derechos respecto de una misma cosa, determina a priori hacia quién se inclinará la protección de la ley.

\subsection{El principio nemo plus iuris}

\textbf{Art.~399 CCyC:} nadie puede transmitir a otro un derecho mejor ni más extenso que el que tiene, salvo las excepciones expresamente establecidas.

La expresión «salvo» abre la puerta a las excepciones. Una de ellas es la \textbf{convalidación} (art.~1885 CCyC) en materia de muebles no registrables.

\textbf{Art.~1885 CCyC (convalidación):} si alguien transmite una cosa de la que no era propietario, y luego la adquiere antes de que se declare la nulidad del acto, la transmisión se convalida.

\subsection{El sub adquirente de buena fe a título oneroso de inmuebles (art.~392 CCyC)}

\textbf{Art.~392 CCyC:} todos los derechos reales o personales transmitidos a terceros sobre un inmueble o muebles registrables por una persona que ha resultado adquirente en virtud de un acto nulo quedan sin ningún valor, pudiendo serle opuestos por quienes tengan interés legítimo, excepto respecto de los sub adquirentes de buena fe y a título oneroso. Los sub adquirentes no pueden ampararse en su buena fe si el acto se ha realizado sin intervención del titular del derecho.

\begin{example}{Caso práctico central}
Diego (enajenante) vende a Miriam (adquirente). Miriam luego vende a Juan (sub adquirente). Después se declara nulo el acto entre Diego y Miriam (por ejemplo, por incapacidad de Diego). ¿Qué sucede con los derechos de Juan?
\end{example}

\begin{table}[H]
\centering
\caption{Seguridad estática y seguridad dinámica}
\small
\begin{tabularx}{\textwidth}{
    >{\raggedright\arraybackslash}p{0.2\textwidth}
    >{\raggedright\arraybackslash}X
    >{\raggedright\arraybackslash}X
}
\toprule
\textbf{Tipo de seguridad} & \textbf{A quién protege} & \textbf{En el ejemplo} \\
\midrule
Estática & Al propietario original (quien tiene el derecho real) & Protege a Diego: el inmueble volvería a él \\
Dinámica & Al adquirente en el tráfico (quien confió en la apariencia) & Protege a Juan: consolida su adquisición \\
Opción del art.~392 CCyC & Seguridad dinámica (el sub adquirente de buena fe y título oneroso) & Juan hace suyo el inmueble \\
\bottomrule
\end{tabularx}
\end{table}

Para que Juan quede protegido se requiere:

\begin{itemize}
    \item Haber adquirido de \textbf{buena fe}.
    \item Haber adquirido a \textbf{título oneroso}.
    \item Que el enajenante (Diego) haya \textbf{participado} de alguna manera en el acto que después fue declarado nulo.
\end{itemize}

\begin{important}
\textbf{Excepción a la excepción:} si hubo sustitución de personalidad (alguien se hizo pasar por Diego), esta protección \textbf{no} aplica a Juan, aunque haya actuado de buena fe y a título oneroso. El enajenante debe haber participado en el acto nulo.
\end{important}

En ningún caso existe protección para Miriam (el adquirente intermedio): ella siempre sufrirá los efectos de la nulidad.

\begin{note}
En el ejercicio profesional se actúa sabiendo de antemano que, en infinidad de situaciones, no es posible atender la justicia de cada uno. Ante la tensión entre distintos derechos, todos atendibles ---como en el caso de Diego, Miriam y Juan---, la norma del ordenamiento jurídico se ve forzada a decidir a priori sobre quién recaerá la protección legal, con la mirada puesta siempre en el orden público y el interés de la comunidad.
\end{note}

\section{Cuadro Cornell: adquisición del dominio}

\begin{longtable}{
    >{\raggedright\arraybackslash}p{0.26\textwidth}
    >{\raggedright\arraybackslash}p{0.66\textwidth}
}
\toprule
\textbf{Preguntas / palabras clave} & \textbf{Notas / respuestas} \\
\midrule
\endfirsthead
\toprule
\textbf{Preguntas / palabras clave} & \textbf{Notas / respuestas} \\
\midrule
\endhead
\midrule
\multicolumn{2}{r}{\small\itshape continúa en la página siguiente} \\
\endfoot
\bottomrule
\endlastfoot

\textbf{¿Qué son los modos de adquisición del dominio?} &
Los hechos y actos jurídicos que la ley prevé y admite para adjudicar en propiedad una cosa mueble o un inmueble a una persona. Los hechos ilícitos (hurto, robo) no son modos reconocidos. \\

\textbf{¿Cómo se clasifican los modos de adquisición?} &
Por el acto (entre vivos o mortis causa), por el objeto (a título universal o singular) y por la causa (originario o derivado). \\

\textbf{¿Qué es la teoría del título y modo?} &
Para adquirir el dominio por modo derivado se necesitan dos elementos: título suficiente (acto jurídico válido que justifica la adquisición) y modo (tradición: entrega efectiva de la cosa). Uno sin el otro no es suficiente. \\

\textbf{¿Qué es el nemo plus iuris y cuál es su excepción?} &
Art.~399: nadie puede transmitir más derechos de los que tiene. Excepción: art.~392, el sub adquirente de buena fe y a título oneroso de un inmueble queda protegido aunque el acto previo fuera nulo, siempre que el enajenante haya participado en ese acto. \\

\textbf{¿Qué es la seguridad estática y la dinámica?} &
Estática: protege al propietario original y su derecho real. Dinámica: protege al tercero que confió en la apariencia jurídica del mercado. El art.~392 opta por la seguridad dinámica para el sub adquirente de buena fe a título oneroso. \\

\textbf{¿Cuándo no protege el art.~392 al sub adquirente?} &
Cuando el acto se realizó sin intervención del titular del derecho, por ejemplo en caso de sustitución de personalidad. El adquirente intermedio nunca queda protegido. \\

\midrule
\multicolumn{2}{p{0.92\textwidth}}{
\textbf{Resumen:}
Para adquirir el dominio por vía derivada rige la teoría del título y modo, pero el sistema admite una importante excepción al principio nemo plus iuris: el sub adquirente de buena fe y a título oneroso de inmuebles queda protegido ante la nulidad del acto anterior, con la condición de que el enajenante haya participado en dicho acto. Esta solución privilegia la seguridad dinámica del tráfico.
} \\
\end{longtable}

% ============================================================
\chapter[Límites al dominio]{Límites al dominio}
% ============================================================

\section{El dominio absoluto: ¿es ilimitado?}

Si el dominio es el derecho más extenso, cabe preguntarse si es un derecho ilimitado, es decir, si el titular puede ejercer las facultades que la ley le confiere del modo en que desee hacerlo.

El ejercicio del derecho de dominio, tanto en Argentina como en la gran mayoría de los sistemas de derecho occidentales, reconoce infinidad de límites, tanto en el interés de los vecinos como en el interés de la comunidad. Estos límites se han ido ampliando a través de los años y continúan modelándose en armonía con el reconocimiento de los derechos humanos, que constituyen el nuevo eje rector de coherencia de todo el ordenamiento nacional.

\section{Evolución histórica del concepto}

\subsection{El Código de Vélez}

El art.~2513 del Código de Vélez Sarsfield otorgaba al titular la facultad de desnaturalizar la cosa y degradarla. Sin embargo, en una nota al art.~2506, el propio Vélez expresaba su pensamiento citando una ley de Partidas: el dominio debe ser ejercido «según Dios y según fuero».

\subsection{La reforma de la Ley 17.711}

La facultad de destruir la cosa fue suprimida en la reforma del Código Civil realizada por la Ley 17.711. Esta reforma incorporó al ordenamiento nacional la noción del \textbf{abuso del derecho}, que fue sumamente relevante en el desarrollo doctrinario posterior.

\begin{table}[H]
\centering
\caption{Evolución del abuso del derecho}
\small
\begin{tabularx}{\textwidth}{
    >{\raggedright\arraybackslash}p{0.22\textwidth}
    >{\raggedright\arraybackslash}X
    >{\raggedright\arraybackslash}X
}
\toprule
\textbf{Momento} & \textbf{Regulación} & \textbf{Criterio de abuso} \\
\midrule
Ley 17.711 & Introdujo la noción de abuso del derecho & Cuando no se tenían en cuenta las razones del legislador al dictar la norma \\
CCyC vigente (art.~10) & Modela todo el ordenamiento jurídico & Cuando el ejercicio contrariara los principios en que está asentado el sistema normativo \\
\bottomrule
\end{tabularx}
\end{table}

\section{La función social de la propiedad}

\textbf{Pacto de San José de Costa Rica, art.~21:} reconoce expresamente el derecho de toda persona al uso y goce de sus bienes, y la facultad de la ley de subordinar tal uso y goce al interés social.

El Pacto de San José de Costa Rica fue incorporado a la Constitución Argentina mediante el art.~75, inc.~22. El CCyC vigente refleja en su articulado este instrumento internacional.

La expresión «función social de la propiedad» fue suprimida del art.~15 del CCyC tal como fue votado por el Congreso, aunque sí figuraba en el proyecto original. Se sostiene que la supresión de la norma no modifica ni aminora la función social de la propiedad que rige en el ordenamiento nacional.

La función social de la propiedad, al tiempo que limita el ejercicio absoluto del derecho del titular dominial, también lo protege del ejercicio ilimitado o abusivo de los demás respecto de él. Forma parte de 14 constituciones provinciales argentinas en forma expresa y es plenamente vigente en el derecho nacional.

\begin{important}
El derecho de dominio en Argentina es un derecho conformado e integrado por la función social de la propiedad. No hay contradicción entre ser propietario y tener límites.
\end{important}

\section{Clasificación de los límites al dominio}

\textbf{Art.~1970 CCyC:} las limitaciones impuestas al dominio privado en el interés público son regidas por el derecho administrativo. El aprovechamiento y uso del dominio sobre inmuebles debe ejercerse de conformidad con las normas administrativas aplicables en cada jurisdicción. Los límites impuestos al dominio en este Capítulo en materia de relaciones de vecindad rigen en subsidio de las normas administrativas aplicables en cada jurisdicción.

El art.~1970 establece dos tipos de restricciones:

\begin{table}[H]
\centering
\caption{Tipos de restricciones al dominio}
\small
\begin{tabularx}{\textwidth}{
    >{\raggedright\arraybackslash}p{0.15\textwidth}
    >{\raggedright\arraybackslash}p{0.19\textwidth}
    >{\raggedright\arraybackslash}X
    >{\raggedright\arraybackslash}p{0.2\textwidth}
}
\toprule
\textbf{Tipo} & \textbf{Fundamento} & \textbf{Regulación} & \textbf{Genera indemnización} \\
\midrule
Interés público & Beneficio de la comunidad en general & Derecho administrativo (ilimitadas en número y clase) & No (en principio) \\
Interés de los vecinos & Convivencia y tutela de la vecindad & Código Civil y Comercial, arts.~1970 a 1982 & No (en principio, salvo daño efectivo) \\
\bottomrule
\end{tabularx}
\end{table}

Son ejemplos de restricciones de interés público: normas de higiene pública, tranquilidad, horarios de actividades, velocidades máximas y mínimas y prohibición de estacionar. Estas limitaciones las sufre el titular en beneficio de la comunidad y no generan indemnización.

\section{Restricciones en interés de los vecinos}

Las restricciones en interés de los vecinos (arts.~1970 a 1982 CCyC) no son taxativas. Las reguladas expresamente en el CCyC son:

\begin{itemize}
    \item Cláusulas de inenajenabilidad (art.~1972).
    \item Inmisiones inmateriales (art.~1973).
    \item Camino de sirga (art.~1974).
    \item Restricciones en materia de aguas (arts.~1975 a 1977).
    \item Luces y vistas (arts.~1978 y 1979).
    \item Árboles, arbustos y otras plantas (art.~1980).
    \item Obligación de soportar instalaciones provisorias y paso de personas en obras linderas.
\end{itemize}

En todos los casos primarán las normas administrativas del lugar si hubiera alguna que modifique las alturas de ventanas, luces, etc. El derecho privado cede ante el orden público.

\subsection{Cláusulas de inenajenabilidad}

\textbf{Art.~1972 CCyC:} las cláusulas de inenajenabilidad son nulas si obligan al propietario a no enajenar a título oneroso o a no constituir derechos reales. Son válidas solo cuando se obligue a no enajenar a persona determinada. Los donantes pueden establecer una prohibición de enajenar por un plazo máximo de 10 años. Los condóminos pueden obligarse a no dividir por un plazo máximo de 10 años.

La razón de esta norma es clara: si el dueño no desea enajenar, no enajena, pues el dominio es perpetuo. Lo que la ley declara nulo es cualquier compromiso a futuro de no enajenar el inmueble a título oneroso o de no constituir derechos reales. Esto refleja la fortaleza del derecho real en pos de la seguridad jurídica.

\subsection{Inmisiones inmateriales}

\textbf{Art.~1973 CCyC:} las molestias que ocasionan el humo, calor, olores, luminosidad, ruidos, vibraciones o inmisiones similares por el ejercicio de actividades en inmuebles vecinos no deben exceder la normal tolerancia teniendo en cuenta las condiciones del lugar y aunque medie autorización administrativa para aquellas. Según las circunstancias del caso, los jueces pueden disponer la cesación de tales molestias o su disminución hasta lo aceptable, y/o la indemnización de los daños. Para disponer el resarcimiento el juez debe ponderar especialmente el respeto debido al uso regular de la propiedad, la prioridad en el uso, el interés general y las exigencias de la producción.

Los tipos de inmisiones son: humo, calor, olores, luminosidad, ruidos y vibraciones. El titular de dominio debe soportarlas si no exceden la normal tolerancia, pero si la exceden, puede reclamar. Los jueces pueden (a) disponer la cesación o disminución de la inmisión, y/o (b) condenar al responsable a indemnizar los daños, aun cuando mediare autorización administrativa.

\subsubsection{El criterio de la normal tolerancia}

La norma provee criterios objetivos que el juez ---y el abogado--- debe ponderar para determinar cuál es la normal tolerancia:

\begin{table}[H]
\centering
\caption{Criterios para determinar la normal tolerancia (art.~1973 CCyC)}
\small
\begin{tabularx}{\textwidth}{
    >{\raggedright\arraybackslash}p{0.2\textwidth}
    >{\raggedright\arraybackslash}X
    >{\raggedright\arraybackslash}X
}
\toprule
\textbf{Criterio} & \textbf{Descripción} & \textbf{Ejemplo} \\
\midrule
Uso regular de la propiedad & Cómo se está usando la propiedad que genera la inmisión & Una fundición que genera calor realiza un uso regular de su propiedad industrial \\
Prioridad en el uso & Quién estaba primero en el lugar & Si la escuela estaba antes que el inmueble adquirido al lado, quien compra sabe a qué se expone \\
Interés general & Si la actividad sirve a la comunidad & Una escuela o un hospital tienen un interés público que pesa en la ponderación \\
Exigencias de la producción & Si la actividad no puede realizarse sin la inmisión & Una fundición no puede funcionar sin generar calor \\
\bottomrule
\end{tabularx}
\end{table}

\begin{example}{Caso real de una fundición}
Una fundición de hierro generaba calor importante que se percibía desde el inmueble contiguo. Al analizar la demanda, el juez ponderó la prioridad en el uso (la fundición estaba primero), el interés general de la producción, las exigencias propias de la actividad (no puede producir sin calor) y el uso regular de la propiedad. La resolución impuso medidas para amortiguar el calor (aislantes y materiales), pero no ordenó el cese de la actividad.
\end{example}

\begin{example}{La escuela}
Quien adquiere un inmueble al lado de una escuela primaria no puede reclamar por el ruido de los niños en el recreo. La prioridad en el uso (la escuela estaba antes) y el interés general (el carácter educativo) son parámetros objetivos de la normal tolerancia. Además, los valores inmobiliarios reflejan la zona: un inmueble junto a una escuela no vale lo mismo que uno en zona residencial.
\end{example}

\subsection{Camino de sirga}

\textbf{Art.~1974 CCyC:} el dueño de un inmueble colindante con cualquiera de las orillas de los cauces o situado en las márgenes de los canales debe dejar libre una franja de quince metros, a la que puede acceder para posibilitar el uso de estos.

El camino de sirga es una franja de 15 metros que debe dejarse libre en los inmuebles limítrofes a ríos o canales. Originalmente la franja era mayor porque se usaba para jalar embarcaciones por la orilla; hoy fue reducida a 15 metros.

% TODO: contenido incompleto o ambiguo en las notas originales (el camino de sirga se describe como restricción de interés público dentro del grupo de restricciones en interés de los vecinos).
\begin{itemize}
    \item Es una restricción de fuente legal de interés público (facilitar la navegación).
    \item No genera indemnización a favor del propietario.
    \item El propietario no puede construir, plantar ni realizar ninguna obra que afecte la navegación o el tránsito de personas o animales en esa franja.
    \item Cualquiera que se sienta perjudicado por obras en esa franja está legitimado para pedir la remoción.
\end{itemize}

\begin{important}
El Estado puede reglamentar ese uso pero no puede construir allí. Para hacerlo, tendría que expropiar, fundado en una ley especial. El titular de dominio tampoco puede obstaculizar el curso de las aguas ni realizar obras que modifiquen su velocidad o dirección (salvo defensivas). Si lo hace, puede ser intimado a removerlas y obligado a indemnizar.
\end{important}

\subsection{Restricciones en materia de aguas}

\textbf{Art.~1976 CCyC:} el titular del fundo inferior debe recibir las aguas, arena y piedras que descienden naturalmente del fundo superior. No puede realizar ninguna obra que altere ese escurrimiento natural. Si lo hiciere, el titular del fundo superior puede pedir la remoción de la obra e, inclusive, la indemnización.

El escurrimiento natural del agua de lluvia, arena y piedras que descienden de inmuebles vecinos debe ser soportado por el titular del fundo inferior.

\subsection{Luces y vistas}

\textbf{Art.~1978 CCyC (vistas):} no pueden tenerse ventanas de vistas que permitan la visión sobre el inmueble lindero, si no se guarda la distancia de tres metros para las ventanas frontales y de sesenta centímetros para las laterales u oblicuas.

\textbf{Art.~1979 CCyC (luces):} las aberturas para luces (que no permiten la vista del vecino sino solo el paso de luz y aire) pueden tenerse en los muros linderos pero deben ser superiores a 1,80 metros desde el piso, salvo que sean fijas o de material no transparente.

\begin{table}[H]
\centering
\caption{Diferencia entre luces y vistas}
\small
\begin{tabularx}{\textwidth}{
    >{\raggedright\arraybackslash}p{0.22\textwidth}
    >{\raggedright\arraybackslash}p{0.22\textwidth}
    >{\raggedright\arraybackslash}X
    >{\raggedright\arraybackslash}p{0.11\textwidth}
}
\toprule
\textbf{Tipo} & \textbf{Función} & \textbf{Distancias exigidas} & \textbf{Referencia} \\
\midrule
Vistas (ventanas frontales) & Permiten ver el inmueble vecino & Mínimo 3 metros del lindero & Art.~1978 \\
Vistas (ventanas laterales u oblicuas) & Permiten ver el inmueble vecino en ángulo & Mínimo 60 cm del lindero & Art.~1978 \\
Luces & Solo permiten la entrada de luz y aire, no la vista & Superiores a 1,80 m desde el piso (salvo fijas o no transparentes) & Art.~1979 \\
\bottomrule
\end{tabularx}
\end{table}

\begin{example}{Conflicto entre vista y elevación}
Si se tiene una ventana que da a la vista, ¿puede impedirse que el vecino eleve su construcción y prive de esa vista y de luz? En principio no puede impedirse que el lindero ejerza su derecho a elevar su construcción. Se menciona un caso en la Patagonia en el que la elevación realizada por un vecino privó directamente la vista del Nahuel Huapi, situación que resulta sumamente interesante en la práctica permanente.
\end{example}

\subsection{Árboles, arbustos y plantas}

Los árboles y plantas generan importantes conflictos de vecindad; se mencionan casos tristemente célebres de extrema violencia, incluyendo homicidios, generados por conflictos de árboles limítrofes.

El CCyC introduce un criterio subjetivo: ya no se fijan distancias mínimas de plantación, sino que nadie puede tener árboles que perjudiquen a los vecinos a ninguna distancia.

\textbf{Art.~1980 CCyC:} nadie puede tener árboles, arbustos u otras plantas que causen molestias que excedan de la normal tolerancia. El vecino perjudicado puede exigir que sean retirados, cortados o podados. Además, puede cortar él mismo las raíces que penetren en su inmueble.

El Código autoriza al vecino a cortar las raíces que penetren en su inmueble, pero \textbf{no} las ramas. Las ramas solo pueden ser cortadas por el propietario o a su pedido.

\begin{example}{El problema de los gomeros}
Las raíces de ciertos árboles, como el gomero, pueden levantar los pisos y las veredas de la casa vecina. El vecino puede cortarlas él mismo, pero respecto de las ramas deberá actuar el propietario.
\end{example}

\section{Regla general final}

\begin{important}
En todos los casos de límites al dominio primarán las normas administrativas del lugar. Si existe una norma administrativa que modifica las alturas de ventanas, luces, distancias de árboles, etc., ella será aplicable por preeminencia del orden público sobre el derecho privado.
\end{important}

Los límites al dominio impuestos por la ley no son taxativos. El orden público es dinámico porque son dinámicos los intereses de la comunidad. Todo hace prever que los límites pueden modificarse en la medida en que respondan a un interés de orden público.

\section{Cuadro Cornell: límites al dominio}

\begin{longtable}{
    >{\raggedright\arraybackslash}p{0.26\textwidth}
    >{\raggedright\arraybackslash}p{0.66\textwidth}
}
\toprule
\textbf{Preguntas / palabras clave} & \textbf{Notas / respuestas} \\
\midrule
\endfirsthead
\toprule
\textbf{Preguntas / palabras clave} & \textbf{Notas / respuestas} \\
\midrule
\endhead
\midrule
\multicolumn{2}{r}{\small\itshape continúa en la página siguiente} \\
\endfoot
\bottomrule
\endlastfoot

\textbf{¿Qué es la función social de la propiedad?} &
La facultad de la ley de subordinar el uso y goce de los bienes al interés social (Pacto de San José, art.~21, incorporado por el art.~75 inc.~22 CN). Aunque el término no figura en el art.~15 del CCyC definitivo, se sostiene que el concepto rige plenamente en el ordenamiento argentino. \\

\textbf{¿Cuáles son los dos tipos de restricciones al dominio?} &
(1) De interés público: reguladas por el derecho administrativo, ilimitadas, sin indemnización. (2) De interés de los vecinos: reguladas por los arts.~1970 a 1982 CCyC, no taxativas, sin indemnización (salvo daño concreto). \\

\textbf{¿Qué es la cláusula de inenajenabilidad y cuándo es válida? (art.~1972)} &
Es nula la cláusula que obliga a no enajenar a título oneroso o a no constituir derechos reales. Es válida solo si obliga a no enajenar a persona determinada. Los donantes pueden prohibir enajenar por un máximo de 10 años y los condóminos pueden obligarse a no dividir por un máximo de 10 años. \\

\textbf{¿Qué son las inmisiones inmateriales?} &
Molestias como humo, calor, olores, luminosidad, ruidos y vibraciones provenientes de inmuebles vecinos. El titular debe soportarlas dentro de la «normal tolerancia». Cuando la exceden, el juez puede ordenar la cesación o disminución y/o la indemnización, aun con autorización administrativa. \\

\textbf{¿Cuáles son los criterios de la «normal tolerancia»?} &
El art.~1973 provee cuatro criterios objetivos: (1) uso regular de la propiedad, (2) prioridad en el uso, (3) interés general, (4) exigencias de la producción. La prioridad en el uso es el criterio más objetivable: quien llega después sabe a qué se expone. \\

\textbf{¿Qué es el camino de sirga?} &
Art.~1974: franja de 15 metros que deben dejar libre los propietarios de inmuebles limítrofes a ríos o canales, para no obstaculizar la navegación. No genera indemnización. El Estado puede reglamentar ese uso pero no puede construir allí sin expropiar. \\

\textbf{¿Qué establece el art.~1976 sobre las aguas?} &
El titular del fundo inferior debe recibir las aguas, arena y piedras que descienden naturalmente del fundo superior y no puede alterar ese escurrimiento con obras; de hacerlo, el titular del fundo superior puede pedir la remoción de la obra e incluso la indemnización. \\

\textbf{¿Cuál es la diferencia entre «luces» y «vistas»?} &
Vistas: aberturas que permiten ver el inmueble vecino; deben estar a un mínimo de 3 metros del lindero (frontales) o 60 cm (laterales u oblicuas). Luces: aberturas que solo permiten el paso de luz y aire, no la vista; deben estar a más de 1,80 m del piso. \\

\textbf{¿Cuál es la regla sobre árboles entre vecinos?} &
Art.~1980: nadie puede tener árboles que causen molestias que excedan la normal tolerancia, a ninguna distancia. El CCyC reemplazó las distancias mínimas por la pauta subjetiva de la normal tolerancia. El vecino puede cortar las raíces que penetren en su inmueble, pero no las ramas. \\

\textbf{¿Qué norma prima en los conflictos de límites al dominio?} &
Las normas administrativas del lugar, por preeminencia del orden público sobre el derecho privado. \\

\midrule
\multicolumn{2}{p{0.92\textwidth}}{
\textbf{Resumen:}
El carácter absoluto del dominio no es ilimitado: el CCyC, en sintonía con el Pacto de San José de Costa Rica (art.~75 inc.~22 CN), reconoce la función social de la propiedad y establece restricciones que el titular debe soportar tanto en interés público como en interés de los vecinos. Estas restricciones ---inmisiones inmateriales, camino de sirga, luces y vistas, árboles--- no generan indemnización en sí mismas y no son taxativas, rigiéndose siempre por las normas administrativas locales cuando estas existieran. El criterio rector en los conflictos de vecindad es la «normal tolerancia», objetivada por la prioridad en el uso, el interés general y las exigencias de la producción.
} \\
\end{longtable}

\end{document}
